\documentclass[11pt]{article}

\usepackage[a4paper,margin=30mm]{geometry}
\usepackage[T1]{fontenc}
\usepackage{lmodern}
\usepackage{microtype}
\usepackage{amsmath,amssymb,amsthm,mathtools}
\usepackage{enumitem}
\usepackage{xcolor}
\usepackage[colorlinks=true,linkcolor=blue!55!black,citecolor=blue!55!black,urlcolor=blue!55!black]{hyperref}

\newtheorem{theorem}{Theorem}[section]
\newtheorem{lemma}[theorem]{Lemma}
\newtheorem{proposition}[theorem]{Proposition}
\newtheorem{corollary}[theorem]{Corollary}
\theoremstyle{remark}
\newtheorem{remark}[theorem]{Remark}

\newcommand{\R}{\mathbb{R}}
\newcommand{\E}{\mathbb{E}}
\newcommand{\SD}{\mathrm{SD}}
\newcommand{\ID}{\mathrm{ID}}
\newcommand{\Log}{\operatorname{Log}}

\title{\bfseries A Jensen Obstruction to Self-Decomposability\\
of the Symmetric \texorpdfstring{$\alpha$}{alpha}-Cauchy Law}
\author{Research note}
\date{August 27, 2026}

\begin{document}
\maketitle

\begin{abstract}
For $\alpha>1$, let $C_\alpha$ have density proportional to
$(1+|x|^\alpha)^{-1}$.  Yano, Yano and Yor asked whether $C_\alpha$
is self-decomposable for $1<\alpha<2$.  Using their Fourier duality
with the symmetric Linnik law, together with the positive Laplace
representation made explicit in Wang's work, we reduce
self-decomposability to positive definiteness of a bounded
background-driving jump function $J_\alpha$.  A gamma scaling,
strict Jensen inequality, and an exact complex beta integral imply
\[
  \int_0^\infty J_\alpha(t)\,dt<0.
\]
Since $J_\alpha$ is integrable, it cannot be a characteristic
function.  Consequently
\[
  C_\alpha\notin\SD,\qquad 1<\alpha<2.
\]
Together with the Cauchy endpoint, this yields
$C_\alpha\in\SD$ if and only if $\alpha=2$ within
$1<\alpha\leq2$.  The argument is entirely analytic; no numerical
sign computation is used.
\end{abstract}

\section{Introduction and provenance}

For $\alpha>1$, the symmetric $\alpha$-Cauchy law is defined by
\begin{equation}\label{eq:density}
 f_\alpha(x)
 =\frac{\alpha\sin(\pi/\alpha)}{2\pi}
  \frac{1}{1+|x|^\alpha},
 \qquad x\in\R.
\end{equation}
Yano, Yano and Yor introduced this family in connection with hitting
times of one-dimensional symmetric stable Levy processes
\cite[Remark 2.9]{YYY2009}.  Their Proposition 2.11 identifies a
Fourier duality between the $\alpha$-Cauchy and $\alpha$-Linnik laws.
Wang subsequently proved
\[
 C_\alpha\in\ID
 \quad\Longleftrightarrow\quad
 1<\alpha\leq2
\]
and supplied a positive Laplace representation of the characteristic
function for $1<\alpha<2$ \cite[Lemma 2.7]{Wang2026}.

The self-decomposability question was not settled in those sources.
The result below is a new deduction from their formulas.

\begin{theorem}\label{thm:main}
For $1<\alpha<2$, the symmetric $\alpha$-Cauchy law is not
self-decomposable:
\[
 C_\alpha\notin\SD.
\]
Consequently, for $1<\alpha\leq2$,
\[
 C_\alpha\in\SD
 \quad\Longleftrightarrow\quad
 \alpha=2.
\]
\end{theorem}

\section{The background-driving reduction}

Let
\[
 \varphi_\alpha(t)=\E[e^{itC_\alpha}],
 \qquad
 \Psi_\alpha(t)=-\log\varphi_\alpha(t).
\]
For $1<\alpha<2$, the Laplace representation recalled in
Section~\ref{sec:mixture} shows that $\varphi_\alpha(t)>0$.

The background-driving Levy process representation of
self-decomposable laws \cite{JurekVervaat1983,Sato1999} gives the
following necessary condition.

\begin{lemma}\label{lem:bdlp}
If $C_\alpha\in\SD$, then the even extension of
\[
 q_\alpha(t)
 :=t\Psi_\alpha'(t)
 =-t\frac{\varphi_\alpha'(t)}{\varphi_\alpha(t)},
 \qquad t>0,
\]
is a continuous negative-definite function.
\end{lemma}

The large-frequency behavior of the Fourier transform of
\eqref{eq:density} is
\begin{equation}\label{eq:phi-asymptotic}
 \varphi_\alpha(t)
 =A_\alpha t^{-\alpha-1}
  \left(1+D_\alpha t^{-\alpha}
  +O(t^{-2\alpha})\right),
\end{equation}
where
\begin{equation}\label{eq:Dalpha}
 D_\alpha
 =-2\cos\left(\frac{\pi\alpha}{2}\right)
   \frac{\Gamma(2\alpha+1)}{\Gamma(\alpha+1)}
 >0.
\end{equation}
The same expansion also follows directly from the Laplace
representation below.  Hence
\begin{equation}\label{eq:q-asymptotic}
 q_\alpha(t)
 =\alpha+1+\alpha D_\alpha t^{-\alpha}
  +O(t^{-2\alpha}),
 \qquad t\longrightarrow\infty.
\end{equation}
In particular, $q_\alpha$ is bounded and
$q_\alpha(t)\to\alpha+1$.

Set
\begin{equation}\label{eq:Jdef}
 J_\alpha(t)
 :=1-\frac{q_\alpha(t)}{\alpha+1}
 =1+\frac{t}{\alpha+1}
      \frac{\varphi_\alpha'(t)}{\varphi_\alpha(t)}.
\end{equation}

\begin{lemma}\label{lem:jumpcf}
If $C_\alpha\in\SD$, then $J_\alpha$ is a characteristic function.
\end{lemma}

\begin{proof}
By Lemma~\ref{lem:bdlp}, $q_\alpha$ is a bounded symmetric
Levy exponent.  Its Levy--Khintchine representation therefore has
the compound-Poisson form
\[
 q_\alpha(t)
 =\int_{\R}(1-\cos(tx))\,\nu_\alpha(dx)
\]
for a finite symmetric measure $\nu_\alpha$ on
$\R\setminus\{0\}$.  Averaging over $0\leq t\leq T$ and using
dominated convergence gives
\[
 \lim_{T\to\infty}\frac1T\int_0^Tq_\alpha(t)\,dt
 =\nu_\alpha(\R).
\]
On the other hand, \eqref{eq:q-asymptotic} shows that the same limit
is $\alpha+1$.  Thus $\nu_\alpha(\R)=\alpha+1$, and
\[
 J_\alpha(t)
 =\frac1{\alpha+1}\int_{\R}\cos(tx)\,\nu_\alpha(dx)
\]
is the characteristic function of the normalized jump measure.
\end{proof}

\section{The positive Laplace representation}\label{sec:mixture}

Put
\[
 \theta=\frac{\pi\alpha}{2},
 \qquad
 p=\alpha+1,
\]
and define
\begin{equation}\label{eq:H}
 H_\alpha(y)
 :=\frac{1}
 {1+2\cos\theta\,y^\alpha+y^{2\alpha}},
 \qquad y>0.
\end{equation}
Since $\theta\in(\pi/2,\pi)$, the denominator in
\eqref{eq:H} is strictly positive.

The Wang--Yano--Yano--Yor representation can be normalized as
\begin{equation}\label{eq:mixture}
 \varphi_\alpha(t)
 =\kappa_\alpha\int_0^\infty
  e^{-ty}y^\alpha H_\alpha(y)\,dy,
 \qquad t\geq0,
\end{equation}
where $\kappa_\alpha>0$ is determined by
$\varphi_\alpha(0)=1$.  Define
\[
 A_\alpha=\kappa_\alpha\Gamma(p),
 \qquad
 F_\alpha(t)
 :=\frac{t^p\varphi_\alpha(t)}{A_\alpha}.
\]
Changing variables $z=ty$ in \eqref{eq:mixture} yields
\begin{equation}\label{eq:gamma-scaling}
 F_\alpha(t)
 =\frac1{\Gamma(p)}
  \int_0^\infty e^{-z}z^{p-1}
      H_\alpha(z/t)\,dz
 =\E[H_\alpha(Z/t)],
\end{equation}
where $Z\sim\Gamma_p$ has unit scale.

Differentiating the definition of $F_\alpha$ gives the exact identity
\begin{equation}\label{eq:J-logF}
 J_\alpha(t)
 =\frac{t}{p}\frac{d}{dt}\log F_\alpha(t).
\end{equation}

\begin{lemma}\label{lem:integrability}
For $1<\alpha<2$, one has
\[
 J_\alpha\in L^1(0,\infty),
 \qquad
 \log F_\alpha\in L^1(0,\infty),
\]
and
\[
 \lim_{t\downarrow0}t\log F_\alpha(t)
 =\lim_{t\to\infty}t\log F_\alpha(t)=0.
\]
Consequently,
\begin{equation}\label{eq:intJ}
 \int_0^\infty J_\alpha(t)\,dt
 =-\frac1p\int_0^\infty\log F_\alpha(t)\,dt.
\end{equation}
\end{lemma}

\begin{proof}
As $y\downarrow0$,
\[
 H_\alpha(y)
 =1-2\cos\theta\,y^\alpha+O(y^{2\alpha}).
\]
Applying this expansion in \eqref{eq:gamma-scaling} gives
\[
 F_\alpha(t)
 =1+D_\alpha t^{-\alpha}+O(t^{-2\alpha}),
 \qquad t\to\infty,
\]
with $D_\alpha$ as in \eqref{eq:Dalpha}.  Hence
\[
 \log F_\alpha(t)=O(t^{-\alpha}),
 \qquad
 J_\alpha(t)=O(t^{-\alpha}).
\]

At the other endpoint, the density in \eqref{eq:mixture} is
asymptotic to $\kappa_\alpha y^{-\alpha}$ as $y\to\infty$.
The standard Laplace asymptotic, or a direct splitting of the
integral at $y=1/t$, gives
\[
 -\varphi_\alpha'(t)
 \sim\kappa_\alpha\Gamma(2-\alpha)t^{\alpha-2},
 \qquad t\downarrow0.
\]
Since $\varphi_\alpha(t)\to1$, it follows that
$q_\alpha(t)=O(t^{\alpha-1})$ and $J_\alpha(t)\to1$.
Furthermore,
\[
 \log F_\alpha(t)=p\log t+O(1),
 \qquad t\downarrow0.
\]
All asserted integrability and boundary statements now follow from
$1<\alpha<2$.  Finally, integrate
\eqref{eq:J-logF} by parts.
\end{proof}

\section{The logarithmic cancellation}

The following elementary complex beta integral is the source of the
exact cancellation.

\begin{lemma}\label{lem:complex-beta}
Let $1<\alpha<2$ and let $a\in\mathbb{C}\setminus(-\infty,0]$,
with all powers and logarithms taken on their principal branches.
Then
\begin{equation}\label{eq:Ialpha}
 I_\alpha(a)
 :=\int_0^\infty
   \frac{\Log(1+a y^\alpha)}{y^2}\,dy
 =\frac{\pi a^{1/\alpha}}{\sin(\pi/\alpha)}.
\end{equation}
\end{lemma}

\begin{proof}
The integral is convergent because its integrand is
$O(y^{\alpha-2})$ at zero and $O(y^{-2}\log y)$ at infinity.
For $a>0$, integration by parts gives
\begin{align*}
 I_\alpha(a)
 &=a\alpha\int_0^\infty
   \frac{y^{\alpha-2}}{1+a y^\alpha}\,dy \\
 &=\frac{\pi a^{1/\alpha}}{\sin(\pi/\alpha)},
\end{align*}
where the last equality is Euler's beta integral.  Both sides are
analytic in $a$ on the slit plane
$\mathbb{C}\setminus(-\infty,0]$; analytic continuation proves
\eqref{eq:Ialpha}.
\end{proof}

\begin{proposition}\label{prop:cancellation}
For $1<\alpha<2$,
\begin{equation}\label{eq:cancellation}
 \int_0^\infty
  \frac{\log H_\alpha(y)}{y^2}\,dy=0.
\end{equation}
\end{proposition}

\begin{proof}
Factor the denominator in \eqref{eq:H}:
\[
 1+2\cos\theta\,y^\alpha+y^{2\alpha}
 =(1+e^{i\theta}y^\alpha)
  (1+e^{-i\theta}y^\alpha).
\]
The two factors are complex conjugates, so their principal logarithms
sum to the real logarithm of the product.  Lemma~\ref{lem:complex-beta}
therefore gives
\begin{align*}
 -\int_0^\infty
  \frac{\log H_\alpha(y)}{y^2}\,dy
 &=I_\alpha(e^{i\theta})+I_\alpha(e^{-i\theta}) \\
 &=\frac{\pi}{\sin(\pi/\alpha)}
   \left(e^{i\theta/\alpha}
         +e^{-i\theta/\alpha}\right).
\end{align*}
Since $\theta/\alpha=\pi/2$, the last parenthesis is
$i-i=0$.
\end{proof}

\section{Strict Jensen obstruction and proof of the theorem}

\begin{proposition}\label{prop:jensen}
For every $1<\alpha<2$,
\[
 \int_0^\infty\log F_\alpha(t)\,dt>0.
\]
\end{proposition}

\begin{proof}
By \eqref{eq:gamma-scaling} and strict concavity of the logarithm,
\begin{equation}\label{eq:strict-jensen}
 \log F_\alpha(t)
 =\log\E[H_\alpha(Z/t)]
 >\E[\log H_\alpha(Z/t)],
 \qquad t>0.
\end{equation}
The inequality is strict because $Z$ is nondegenerate and
$H_\alpha$ is nonconstant.

The required exchange of expectation and integration is absolute.
Indeed,
\[
 \int_0^\infty
  \frac{|\log H_\alpha(y)|}{y^2}\,dy<\infty:
\]
near zero the integrand is $O(y^{\alpha-2})$, and at infinity it is
$O(y^{-2}\log y)$.  With the substitution $y=Z/t$,
\begin{align*}
 \E\left[
  \int_0^\infty|\log H_\alpha(Z/t)|\,dt
 \right]
 &=\E[Z]\int_0^\infty
   \frac{|\log H_\alpha(y)|}{y^2}\,dy \\
 &<\infty.
\end{align*}
Thus Proposition~\ref{prop:cancellation} implies
\[
 \int_0^\infty
  \E[\log H_\alpha(Z/t)]\,dt
 =\E[Z]\int_0^\infty
   \frac{\log H_\alpha(y)}{y^2}\,dy
 =0.
\]
Integrating the strict inequality \eqref{eq:strict-jensen} proves
the claim.
\end{proof}

\begin{proof}[Proof of Theorem~\ref{thm:main}]
Combining Proposition~\ref{prop:jensen} with
\eqref{eq:intJ} yields
\begin{equation}\label{eq:negative-integral}
 \int_0^\infty J_\alpha(t)\,dt<0.
\end{equation}
Suppose for contradiction that $C_\alpha\in\SD$.
Lemma~\ref{lem:jumpcf} says that the even function $J_\alpha$ is a
characteristic function.  By Lemma~\ref{lem:integrability}, it is
integrable.  Fourier inversion would therefore give its probability
density $r_\alpha$, continuous on $\R$, with
\[
 r_\alpha(0)
 =\frac1{2\pi}\int_{\R}J_\alpha(t)\,dt
 =\frac1\pi\int_0^\infty J_\alpha(t)\,dt<0.
\]
This is impossible for a probability density, and hence
$C_\alpha\notin\SD$.

At $\alpha=2$, $C_2$ is the standard Cauchy law, which is stable and
therefore self-decomposable.  The endpoint statement follows.
\end{proof}

\section{The explicit case \texorpdfstring{$\alpha=3/2$}{alpha=3/2}}

For $\alpha=3/2$,
\[
 \theta=\frac{3\pi}{4},
 \qquad p=\frac52,
\]
and the positive mixture becomes
\[
 \varphi_{3/2}(t)
 =\kappa\int_0^\infty e^{-ty}
  \frac{y^{3/2}}
  {1-\sqrt2\,y^{3/2}+y^3}\,dy.
\]
Writing
\[
 H(y)=\frac1{1-\sqrt2\,y^{3/2}+y^3},
 \qquad Z\sim\Gamma_{5/2},
\]
one has
\[
 F(t)=\frac{t^{5/2}\varphi_{3/2}(t)}
 {\kappa\Gamma(5/2)}
 =\E[H(Z/t)]
\]
and
\[
 J_{3/2}(t)=\frac{2t}{5}(\log F(t))'.
\]
The factorization
\[
 1-\sqrt2\,y^{3/2}+y^3
 =(1+e^{3\pi i/4}y^{3/2})
  (1+e^{-3\pi i/4}y^{3/2})
\]
and Lemma~\ref{lem:complex-beta} give
\[
 \int_0^\infty
  \frac{\Log(1+a y^{3/2})}{y^2}\,dy
 =\frac{\pi a^{2/3}}{\sin(2\pi/3)}.
\]
For $a=e^{\pm3\pi i/4}$, the principal powers are
$a^{2/3}=e^{\pm i\pi/2}=\pm i$, and hence
\[
 \int_0^\infty\frac{\log H(y)}{y^2}\,dy=0.
\]
Strict Jensen then yields
\[
 \int_0^\infty J_{3/2}(t)\,dt
 =-\frac25\int_0^\infty\log F(t)\,dt<0,
\]
which is the explicit obstruction proving
$C_{3/2}\notin\SD$.

\begin{remark}
The negative sign above is not inferred from a numerical Fourier
inversion.  It follows solely from strict Jensen inequality and the
exact cancellation in the complex beta integral.
\end{remark}

\begin{thebibliography}{9}

\bibitem{JurekVervaat1983}
Z.~J. Jurek and W. Vervaat,
\newblock An integral representation for self-decomposable Banach
space valued random variables,
\newblock \emph{Z. Wahrscheinlichkeitstheorie verw. Gebiete}
\textbf{62} (1983), 247--262.
\newblock \href{https://doi.org/10.1007/BF00538800}
{doi:10.1007/BF00538800}.

\bibitem{Sato1999}
K.-I. Sato,
\newblock \emph{Levy Processes and Infinitely Divisible
Distributions},
\newblock Cambridge Studies in Advanced Mathematics 68,
Cambridge University Press, 1999.

\bibitem{Wang2026}
M. Wang,
\newblock Infinite divisibility of $\alpha$-Cauchy distributions,
\newblock arXiv:2512.23164v3, 2026; see also the supplied
\emph{Electronic Communications in Probability} page proof.
\newblock \url{https://arxiv.org/abs/2512.23164}.

\bibitem{YYY2009}
K. Yano, Y. Yano and M. Yor,
\newblock On the laws of first hitting times of points for
one-dimensional symmetric stable Levy processes,
\newblock in \emph{Seminaire de Probabilites XLII},
Lecture Notes in Mathematics 1979, Springer, 2009, pp. 187--227.
\newblock \href{https://doi.org/10.1007/978-3-642-01763-6_8}
{doi:10.1007/978-3-642-01763-6\_8};
\url{https://arxiv.org/abs/0811.2046}.

\end{thebibliography}

\end{document}
